\chapter{Carry Across Asset Classes}\label{ch:s1:carry-across-asset-classes}

Borrow in a currency that pays little interest, lend in one that pays a lot, and collect the difference as long as the exchange rate does not move
against you. It usually does not move enough, and the currency carry trade has paid for decades; then, in a few weeks, the high-yielding currencies fall
together and years of income go. Koijen, Moskowitz, Pedersen and Vrugt showed that the same idea, the return an asset earns if its price does not change,
predicts returns in equities, bonds, commodities, credit and options as well as currencies, and that the carry strategies of all classes do badly together
in global recessions. On the synthetic universe a \omterm{def:s1:carry-across-asset-classes:diversified}{diversified carry} book earns a Sharpe ratio of 0.95 when all the carry is earned, and loses 26\% of its
capital, at 10\% volatility, in the second stock crash. The build is \texttt{firm.carrystrat}.

\section{Carry everywhere}

Book~1, chapter~7 defined carry as what a position earns if prices stay where they are. For a futures position it has a common form: with a spot or
reference price $S$ and a futures price $F$ for delivery in $\tau$ years, the carry is $C = (S - F)/(F\tau)$, the futures price's convergence to the spot
if the spot does not move. Each class gives it a name.
\begin{itemize}
\item \textbf{Currencies}: the forward discount, the interest rate of the foreign currency minus the domestic one.
\item \textbf{Equity indices}: the dividend yield minus the financing rate.
\item \textbf{Government bonds}: the yield over the financing rate plus the roll-down, the price gain as a bond ages down an upward-sloping curve.
\item \textbf{Commodities}: the roll yield, positive in backwardation (futures below spot) and negative in contango, set by storage costs and the
convenience yield (Book~3, chapter~10).
\end{itemize}
Koijen and co-authors write a futures position's expected excess return as its carry plus the expected change in the spot price. If spot prices do not
move on average, the whole carry is earned; if they moved to cancel it, as the expectations hypothesis and uncovered interest parity would have them do,
none would be. \texttt{firm.synthfut} plants a share $\lambda$ of the carry as expected return: $\lambda = 0.5$ in the universe of all these chapters, and
$\lambda = 1$ in a variant where prices do not move against carry.

\begin{definition}[Carry strategy]\label{def:s1:carry-across-asset-classes:carry}
A \emph{carry strategy}\index{carry strategy} holds long positions in the assets of a class with the highest carry and short positions in those with
the lowest, with weights that increase with the carry's rank, rebalanced as carries change; it earns the carry spread and bears the risk that prices
move against it.
\end{definition}

In the synthetic universe carry is read from the curve (\cref{lst:s1:carry-across-asset-classes:measures}): the log slope between the first two monthly
contracts, 21 trading days apart, annualised. For most markets this equals the planted carry exactly, because the model's curves are straight lines in
maturity. For the three seasonal commodities it does not: their curves carry the annual cycle of their spot prices, and the one-month slope mostly
measures where the season is going. Its correlation with the true carry is 0.05. The slope between contracts twelve months apart cancels the cycle and
has a correlation of 1.00. The chapter's books use it for commodities.

\section{Carry books by class}

Each class's book weights its ten markets by the centred rank of their carry, with a gross exposure of one, rebalanced daily at 2, 1, 2 and 4 basis
points per unit traded by class (\cref{lst:s1:carry-across-asset-classes:books}); each is scaled over the whole sample to 10\% volatility, and the
diversified book gives each class the same weight.

\begin{definition}[Diversified carry]\label{def:s1:carry-across-asset-classes:diversified}
\emph{Diversified carry}\index{diversified carry} is a portfolio of carry strategies in several asset classes, each scaled to a similar risk, which
earns the carry premia of all classes at a volatility lower than any one, as long as their returns are not correlated.
\end{definition}

\begin{center}
\small
\begin{tabular}{@{}lccccc@{}}
\toprule
Sharpe ratio after costs & equities & bonds & currencies & commodities & diversified\\
\midrule
half of the carry earned ($\lambda = 0.5$) & 0.17 & 0.22 & 0.03 & 0.27 & 0.34\\
all of the carry earned ($\lambda = 1$) & 0.28 & 0.53 & 0.49 & 0.59 & 0.95\\
\midrule
monthly skewness ($\lambda = 1$) & 0.13 & $-0.09$ & $-0.36$ & $-0.16$ & $-0.20$\\
\bottomrule
\end{tabular}
\end{center}

Two things set the numbers. The share of carry earned doubles or triples every class's Sharpe ratio (\cref{fig:s1:carry-across-asset-classes:classes}),
which is why the empirical question of whether prices move against carry is the whole question. And the diversified book beats every class: four books
with little correlation in normal times, each at 10\% volatility, combine into one with a Sharpe ratio close to twice that of the average class. The
currency book at $\lambda = 0.5$ earns almost nothing because its crashes, described next, take back what its carry pays.

\begin{omfigure}
\begin{tikzpicture}
\begin{axis}[width=10.5cm,height=5.2cm,ybar,bar width=11pt,ylabel={\small Sharpe ratio after costs},tick label style={font=\footnotesize},
  symbolic x coords={equity,bond,currency,commodity,diversified},xtick=data,ymin=0,ymax=1.1,enlarge x limits=0.14,grid=major,
  grid style={gray!25},legend columns=2,legend style={font=\scriptsize,at={(0.5,-0.18)},anchor=north,draw=none,fill=none,
  /tikz/every even column/.append style={column sep=8pt}},area legend,table/col sep=comma]
\addplot[fill=omExo!50,draw=omExo] table[x=book,y=half]{figdata/strategies-1/20-carry-across-asset-classes/classes.csv};
\addplot[fill=omThm!60,draw=omThm] table[x=book,y=full]{figdata/strategies-1/20-carry-across-asset-classes/classes.csv};
\legend{half of the carry earned,all of the carry earned}
\end{axis}
\end{tikzpicture}

\omcaption{Carry books on the synthetic universe, thirty years: Sharpe ratios after costs by class and for the diversified book, when half of each
market's carry is earned as expected return and when all of it is. Data: \texttt{s1\_carry.summary}.}\label{fig:s1:carry-across-asset-classes:classes}
\end{omfigure}

\section{Carry crashes}

\begin{definition}[Carry crash]\label{def:s1:carry-across-asset-classes:crash}
A \emph{carry crash}\index{carry crash} is a sudden, large loss of a \omterm{def:s1:carry-across-asset-classes:carry}{carry strategy} when high-carry assets fall and low-carry assets rise together,
typically as leveraged carry positions are unwound in a fall of risk appetite or funding liquidity; it makes the strategy's returns negatively skewed.
\end{definition}

Brunnermeier, Nagel and Pedersen documented the pattern in currencies: exchange rate moves between high- and low-interest-rate currencies are negatively
skewed, because carry trades are unwound suddenly when risk appetite and funding liquidity fall. Lustig, Roussanov and Verdelhan found why the premium
exists: high-rate currencies load more on a global risk factor, so the carry investor holds global risk, particularly in bad times. The premium is
compensation for losing when losing hurts most.

\begin{dated}{2026-09}{The August 2024 carry unwind}\label{dat:s1:carry-across-asset-classes:yen}
The Bank for International Settlements described the market turbulence of early August 2024 in its September 2024 Quarterly Review: after Federal
Reserve and Bank of Japan meetings perceived as somewhat hawkish, markets overreacted to a disappointing US labour market release, currency carry trades
unwound amid changing rate expectations and higher volatility, and the funding currencies, predominantly the yen, appreciated sharply for a short time.
The review cites an estimate of hedge funds' foreign exchange forward positions of about \$160 billion.
\end{dated}

The synthetic universe plants the crash: in each stock crash, high-carry currencies lose 8\% per unit of their carry's z-score over the hundred days.
The currency book loses 29.0\%, 24.9\% and 28.0\% in the three crashes ($\lambda = 1$, at 10\% volatility) and has the most negative monthly skewness.
Crashes also hit the other classes, because the crash moves prices in ways that cut across carry. The diversified book loses 12.9\%, 26.3\% and
13.1\%: in the second crash it lost more than any single class, because every class lost at once (\cref{fig:s1:carry-across-asset-classes:paths}). That
is Koijen and co-authors' finding in miniature: the classes are uncorrelated on average and correlated when it matters.

\begin{omfigure}
\begin{tikzpicture}
\begin{axis}[width=11cm,height=5.4cm,xlabel={\small year},ylabel={\small cumulative return (\%)},tick label style={font=\footnotesize},
  xmin=1,xmax=30,ymin=-80,ymax=320,grid=major,grid style={gray!25},legend columns=2,legend style={font=\scriptsize,at={(0.5,-0.3)},
  anchor=north,draw=none,fill=none,/tikz/every even column/.append style={column sep=8pt}},table/col sep=comma]
\fill[omExo!15] (axis cs:7.5,-80) rectangle (axis cs:7.897,320);
\fill[omExo!15] (axis cs:18,-80) rectangle (axis cs:18.397,320);
\fill[omExo!15] (axis cs:27,-80) rectangle (axis cs:27.397,320);
\addplot[omThm,thick] table[x=year,y=diversified]{figdata/strategies-1/20-carry-across-asset-classes/paths.csv};
\addplot[omDef,thick] table[x=year,y=currency]{figdata/strategies-1/20-carry-across-asset-classes/paths.csv};
\legend{diversified carry,currency carry}
\end{axis}
\end{tikzpicture}

\omcaption{The diversified and currency carry books on the synthetic universe when all carry is earned, each at 10\% volatility, as cumulative sums of
daily returns; the shaded bands are the three stock crashes. Data: \texttt{s1\_carry.paths}.}\label{fig:s1:carry-across-asset-classes:paths}
\end{omfigure}

Carry and trend make a natural pair. The \omterm{def:s1:carry-across-asset-classes:diversified}{diversified carry} book's daily returns have a correlation of $-0.00$ with chapter~19's blended trend book
($\lambda = 0.5$), and in the crashes they move in opposite directions: trend gains, carry loses. A book holding both has smoother returns than
either.

\section{Carry on the WTI curve}

The EIA's settlements of the second and third WTI contracts, a month apart, give crude oil's carry from 1985 to 2024 (\cref{fig:s1:carry-across-asset-classes:wti}).
Using them rather than the first contract keeps clear of the contract in its last days before expiry. The average was 1.2\% a year; the curve was in
backwardation on 51.7\% of days. The extremes were $+126\%$ a year on 14 January 1991 and $-575\%$ on 21 April 2020, the day after the May
contract settled below zero (Book~3, chapter~2). A rule that holds the rolled future long in backwardation and short in contango, at a 40\% volatility
target and 2 basis points a trade, earned 24.8\% a year from 1986 to 2024 at a Sharpe ratio of 0.61; this is one market, in sample, with no parameter to
fit beyond the sign.

\begin{omfigure}
\begin{tikzpicture}
\begin{axis}[width=11cm,height=5cm,xlabel={\small year},ylabel={\small carry, contracts 2--3 (\% a year)},tick label style={font=\footnotesize},
  xmin=1985,xmax=2024.5,ymin=-200,ymax=150,grid=major,grid style={gray!25},table/col sep=comma,/pgf/number format/1000 sep={}]
\addplot[omThm] table[x=year,y=carry_pct]{figdata/strategies-1/20-carry-across-asset-classes/wti_carry.csv};
\addplot[black!40,dashed] coordinates {(1985,0) (2024.5,0)};
\end{axis}
\end{tikzpicture}

\omcaption{The carry of NYMEX WTI crude oil futures, $12\ln(F_2/F_3)$, from the EIA's daily settlements of the second and third contracts, every fifth
trading day, 1985--2024; positive in backwardation. The scale stops at $-200\%$: the daily low was $-575\%$ on 21 April 2020. Data:
\texttt{s1\_carry.wti}.}\label{fig:s1:carry-across-asset-classes:wti}
\end{omfigure}

\section{Strategy files}

\begin{strategyfile}{Currency carry}\label{strat:s1:carry-across-asset-classes:fx}
\sfield{Who pays you, and why} Investors who hedge or fund in low-rate currencies, and the market's price for global crash risk.
\sfield{Instruments and venues} Currency forwards and futures of developed and liquid emerging currencies.
\sfield{Signal} The forward discount (interest differential), ranked.
\sfield{Sizing and execution} Long high-rate, short low-rate currencies, volatility-scaled; monthly forwards rolled.
\sfield{Costs} Low; forward points and bid--ask spreads.
\sfield{How it dies} Crashes when leveraged carry positions unwind together.
\sfield{Horizon, capacity, infrastructure} Months; large capacity in the major currencies.
\sfield{Backtest honestly} Forward rates as quoted, not interest-rate proxies; the crash years in the sample.
\sfield{Sources} Brunnermeier, Nagel and Pedersen (2008); Lustig, Roussanov and Verdelhan (2011); the dated box.
\end{strategyfile}

\begin{strategyfile}{Bond carry}\label{strat:s1:carry-across-asset-classes:bond}
\sfield{Who pays you, and why} Investors paying a term premium to hold short duration.
\sfield{Instruments and venues} Government bond futures across countries.
\sfield{Signal} The yield over the financing rate plus roll-down, ranked.
\sfield{Sizing and execution} Duration-neutral or volatility-scaled long--short.
\sfield{Costs} Low.
\sfield{How it dies} Sharp rises in rates at the front of steep curves.
\sfield{Horizon, capacity, infrastructure} Months; a curve model per country.
\sfield{Backtest honestly} Carry from the curve known at the time, with the cheapest-to-deliver of each future.
\sfield{Sources} Koijen, Moskowitz, Pedersen and Vrugt (2018).
\end{strategyfile}

\begin{strategyfile}{Commodity roll yield}\label{strat:s1:carry-across-asset-classes:commodity}
\sfield{Who pays you, and why} Hedgers paying to lay off inventory risk; the market's price for scarcity.
\sfield{Instruments and venues} Commodity futures.
\sfield{Signal} The curve's slope, over twelve months for seasonal commodities.
\sfield{Sizing and execution} Long backwardated, short contango commodities, volatility-scaled.
\sfield{Costs} Moderate; rolls.
\sfield{How it dies} Storage gluts and squeezes that invert the relation between slope and return.
\sfield{Horizon, capacity, infrastructure} Months.
\sfield{Backtest honestly} The contracts that would have been held, not the nearby; seasonal slopes separated from carry.
\sfield{Sources} Koijen, Moskowitz, Pedersen and Vrugt (2018); this chapter's seasonal test.
\end{strategyfile}

\begin{strategyfile}{Diversified carry}\label{strat:s1:carry-across-asset-classes:diversified}
\sfield{Who pays you, and why} The carry premia of all classes at once.
\sfield{Instruments and venues} Futures and forwards in equities, bonds, currencies and commodities.
\sfield{Signal} Carry in each class, ranked within the class.
\sfield{Sizing and execution} Each class at the same risk; rebalanced monthly or as carries change.
\sfield{Costs} Low to moderate.
\sfield{How it dies} Correlated losses in global recessions, when diversification disappears.
\sfield{Horizon, capacity, infrastructure} Months; a data pipeline for every class's carry.
\sfield{Backtest honestly} Carry measures defined before the test; recessions in the sample.
\sfield{Sources} Koijen, Moskowitz, Pedersen and Vrugt (2018); this chapter: 0.95 with all carry earned, a 26.3\% loss in a crash.
\end{strategyfile}

\begin{strategyfile}{WTI curve carry}\label{strat:s1:carry-across-asset-classes:wti}
\sfield{Who pays you, and why} Producers and consumers hedging, who pay for immediacy in backwardation and for storage in contango.
\sfield{Instruments and venues} NYMEX WTI futures, rolled.
\sfield{Signal} The sign of the slope between the second and third contracts.
\sfield{Sizing and execution} Long in backwardation, short in contango, at a volatility target.
\sfield{Costs} Rolls and spreads.
\sfield{How it dies} As one market, it has no diversification; storage limits (April 2020).
\sfield{Horizon, capacity, infrastructure} Weeks to months.
\sfield{Backtest honestly} The rolled series; the slope from contracts clear of expiry.
\sfield{Sources} EIA settlements; this chapter: a Sharpe ratio of 0.61, 1986--2024, in sample.
\end{strategyfile}

\section{Tutorial: autumn in the synthetic market}\label{tut:s1:carry-across-asset-classes:tut}

\begin{tutorial}
\textbf{Goal.} Read carry from the synthetic curves, build carry books by class and diversified under two assumptions about how much carry is earned,
measure them in the crashes, and read crude oil's carry from the EIA curve. \textbf{End state:} the table and the three figures.
\begin{tutsteps}
\item \textbf{Carry measures and rank weights}.
\omcode{code/firm/carrystrat/firm_carrystrat.py}{26}{47}{Carry by class and rank weights within groups.}\label{lst:s1:carry-across-asset-classes:measures}
\item \textbf{The books}: one per class, then the diversified book.
\omcode{code/strategies-1/20-carry-across-asset-classes/python/s1_carry.py}{60}{72}{Carry books by class and diversified.}\label{lst:s1:carry-across-asset-classes:books}
\item \textbf{Run} \texttt{carry\_quality()}, \texttt{summary} with the mixed signal at premiums 0.5 and 1.0, \texttt{trend\_correlation()},
\texttt{wti()}, and \texttt{fig\_carry.py}.
\end{tutsteps}
\textbf{What to change next.} Scale each class's book ex ante with an EWMA forecast instead of in sample; add a crash filter that cuts currency carry
when volatility jumps; combine carry with chapter~19's trend book and measure the combined drawdowns.
\end{tutorial}

\section{Build: carry strategies}\label{bld:s1:carry-across-asset-classes:carrystrat}

\begin{build}
\bfield{Purpose} Carry measures for every class, carry books by rank, and crash diagnostics.
\bfield{Interface} \texttt{curve\_carry(f\_near, f\_far, days)}, \texttt{forward\_discount(log\_spot, log\_fwd, days)}, \texttt{dividend\_carry(div\_yield,
rate)}, \texttt{rank\_weights(x, groups)}, \texttt{book(r, w, cost)}, \texttt{scale(pnl, target, start)}, \texttt{windows(pnl, spans)},
\texttt{skew\_monthly(pnl, start, days)}.
\bfield{Rules} Carry annualised on 252 days; ranks within each group; weights with a gross of one per group; P\&L from weights set at the previous close.
\bfield{Acceptance tests} \texttt{code/firm/carrystrat/tests/}: carry from two prices, a forward and a dividend yield by hand; rank weights by group;
P\&L, costs, windows and a skewed series.
\bfield{Stretch} Ex-ante scaling; carry from full curves with roll-down; crash filters.
\end{build}

\begin{omsources}
\item R.~S.~J.~Koijen, T.~J.~Moskowitz, L.~H.~Pedersen and E.~B.~Vrugt, ``Carry'', \emph{Journal of Financial Economics} 127(2), 2018.
\item M.~K.~Brunnermeier, S.~Nagel and L.~H.~Pedersen, ``Carry trades and currency crashes'', \emph{NBER Macroeconomics Annual} 23, 2008.
\item H.~Lustig, N.~Roussanov and A.~Verdelhan, ``Common risk factors in currency markets'', \emph{Review of Financial Studies} 24(11), 2011.
\item Bank for International Settlements, ``Carry off, carry on'', \emph{BIS Quarterly Review}, September 2024.
\end{omsources}

\section{Exercises}

\begin{exercise}[$\star$]\label{exo:s1:carry-across-asset-classes:1}
Two futures contracts 21 trading days apart are priced at 100 and 98. What is the carry, annualised on 252 days?
\end{exercise}

\begin{exercise}[$\star$]\label{exo:s1:carry-across-asset-classes:2}
The domestic interest rate is 1\% and a foreign currency's is 5\%. What is the carry of a long position in the foreign currency's forward? And of an
equity index future with a dividend yield of 2\% when the financing rate is 4.5\%?
\end{exercise}

\begin{exercise}[$\star$]\label{exo:s1:carry-across-asset-classes:3}
Why is the slope between the first two contracts a poor measure of carry for natural gas?
\end{exercise}

\begin{exercise}[$\star\star$]\label{exo:s1:carry-across-asset-classes:4}
The diversified book runs at 10\% volatility with a Sharpe ratio of 0.95 and loses 26.3\% in a hundred-day crash. How many of its hundred-day standard
deviations is that, and what did it expect to earn over the hundred days?
\end{exercise}

\begin{exercise}[$\star\star$]\label{exo:s1:carry-across-asset-classes:5}
Why does the diversified book lose more than any single class in the second crash?
\end{exercise}

\begin{exercise}[$\star\star$]\label{exo:s1:carry-across-asset-classes:6}
Explain, with Koijen and co-authors' decomposition, why the share of carry earned is ``the whole question''.
\end{exercise}

\begin{exercise}[$\star\star\star$]\label{exo:s1:carry-across-asset-classes:7}
\emph{Coding.} Run \texttt{books} with the signal \texttt{near} and a premium of 1.0, which reads commodity carry from the first two contracts. What happens to the commodity and diversified
Sharpe ratios, and why?
\end{exercise}

\begin{exercise}[$\star\star\star$]\label{exo:s1:carry-across-asset-classes:8}
\emph{Find the flaw.} ``Our currency carry book has had a Sharpe ratio of 1.0 for ten years with a maximum drawdown of 8\%: it is a low-risk strategy.''
\end{exercise}

\section{Problem: Autumn in the Synthetic Market}

\begin{problem}[{Weekend problem --- carry through a crash}]\label{pb:s1:carry-across-asset-classes:1}
The chapter's synthetic universe, the WTI curve and the public record.

\medskip\noindent\textbf{Part I --- Carry.}
\begin{enumerate}
\item Define carry for a futures position and name it in each class.
\item Write Koijen and co-authors' decomposition of expected return.
\item Define a \omterm{def:s1:carry-across-asset-classes:carry}{carry strategy} and \omterm{def:s1:carry-across-asset-classes:diversified}{diversified carry}.
\item How is carry read from the synthetic curves, and what goes wrong for seasonal commodities?
\end{enumerate}

\medskip\noindent\textbf{Part II --- The books.}
\begin{enumerate}[resume]
\item Describe the books: weights, costs, scaling.
\item Give the Sharpe ratios by class and diversified for $\lambda = 0.5$ and $\lambda = 1$.
\item Why does the diversified book beat every class?
\item Why does the currency book earn almost nothing at $\lambda = 0.5$?
\end{enumerate}

\medskip\noindent\textbf{Part III --- Crashes.}
\begin{enumerate}[resume]
\item Define a \omterm{def:s1:carry-across-asset-classes:crash}{carry crash}; what did Brunnermeier, Nagel and Pedersen find?
\item What did Lustig, Roussanov and Verdelhan find?
\item Give the currency and diversified books' losses in the three crashes.
\item What does the dated box describe?
\end{enumerate}

\medskip\noindent\textbf{Part IV --- The verdict.}
\begin{enumerate}[resume]
\item State the \emph{named result}: the \omterm{def:s1:carry-across-asset-classes:diversified}{diversified carry} book's Sharpe ratio and its loss in the crash, against each class alone.
\item How do carry and trend combine?
\item What was WTI's carry from 1985 to 2024, and what did the sign rule earn?
\item What do the extremes of WTI's carry correspond to?
\item How would you backtest a \omterm{def:s1:carry-across-asset-classes:carry}{carry strategy} honestly?
\item Which strategy file is most exposed to storage limits?
\item What did Koijen and co-authors find about carry across classes in recessions?
\item In one sentence: what is a carry investor paid for?
\end{enumerate}
\end{problem}

\section{Interview questions}

\begin{interviewq}[$\star$ \iqroles{researcher}]\label{iq:s1:carry-across-asset-classes:1}
What is carry for a currency, a bond future and a commodity future?
\end{interviewq}

\begin{interviewq}[$\star\star$ \iqroles{researcher}]\label{iq:s1:carry-across-asset-classes:2}
Why would uncovered interest parity imply that currency carry earns nothing, and what does the evidence say?
\end{interviewq}

\begin{interviewq}[$\star\star$ \iqroles{trader}]\label{iq:s1:carry-across-asset-classes:3}
Volatility jumps and the yen rallies 3\% in a day. What do you do with a currency carry book?
\end{interviewq}

\begin{interviewq}[$\star\star$ \iqroles{risk}]\label{iq:s1:carry-across-asset-classes:4}
How would you measure the crash risk of a carry book that has never had a bad month in its history?
\end{interviewq}

\begin{interviewq}[$\star\star$ \iqroles{developer}]\label{iq:s1:carry-across-asset-classes:5}
What data does a \omterm{def:s1:carry-across-asset-classes:diversified}{diversified carry} book need each day, and what can go wrong with it?
\end{interviewq}

\begin{interviewq}[$\star\star\star$ \iqroles{researcher}]\label{iq:s1:carry-across-asset-classes:6}
A futures price is $F = S e^{-c\tau}$ with a constant carry $c$ and a spot $S$ following a martingale. Show that the position's expected return is $c$
per unit time, and find the expected return when instead $E[dS/S] = -(1 - \lambda) c\, dt$.
\end{interviewq}
